\chapter{Problem Formulation}
The goal of this thesis is to estimate the 3D motion of athletes from monocular drone footage.\\ 
Specifically, the task consists in predicting the 3D HPE of an athlete over time, represented as a sequence of 3D joint coordinates.
Given a video acquired from a UAV-mounted camera, the objective is to recover the spatio-temporal configuration of the human body despite challenging acquisition conditions such as motion blur, occlusions, and dynamic camera motion.


\section{Input Representation}
The input to the model is an RGB video sequence :
\begin{equation}
    x \in \mathbb{R}^{T \times H \times W \times C}
\end{equation}
where $T$ is the temporal length of the video, $H$ and $W$ are the spatial dimensions of each frame and $C$ is the number of channels.\\
The video is assumed to be captured by a moving drone-mounted monocular camera, resulting in non-static backgrounds and varying viewpoints.


\section{Output Representation}
The output is a sequence of 3D human poses, one for each video frame. 
Each pose is represented as a set of $J$ 3D joint coordinates:
\begin{equation}
    y \in \mathbb{R}^{T \times J \times P}
\end{equation}
where $T$ is the temporal length of the sequence and $P = 3$ represents for each joint its $(x, y, z)$ coordinates in a root-relative reference frame, obtained by translating all joints such that the root joint is located at 
the origin.


\section{Task Definition}
The problem addressed in this thesis is to learn a function $f_\theta$, parameterized by $\theta$, that maps the input video to a sequence of 3D human poses:
\begin{equation}
    f_\theta: x \in \mathbb{R}^{T \times H \times W \times C} \rightarrow y \in \mathbb{R}^{T \times J \times P}
\end{equation}


\section{Learning Objective}
Given a training set of video sequences and corresponding ground-truth 3D poses, the learning objective is to minimize the discrepancy between predicted and ground-truth joint positions.
This is typically achieved by minimizing the \textit{MPJPE} loss (see \ref{mpjpe}) over all joints and time steps:
\begin{equation}
    \text{MPJPE} = \frac{1}{T} \frac{1}{J} \sum_{t=1}^{T} \sum_{j=1}^{J} \| \hat{\mathbf{y}_{t,j}} - \mathbf{y}_{t,j} \|
\end{equation}


\section{Challenges and Constraints}
Estimating 3D human pose from drone footage presents several challenges:
\begin{itemize}
    \item \textbf{Motion blur and low resolution}, caused by fast camera and subject motion;
    \item \textbf{Self-occlusions and partial visibility} of the human body;
    \item \textbf{Dynamic camera motion}, which prevents the use of static background assumptions;
    \item \textbf{Limited availability of labeled datasets}, as only a few public UAV-based 3D human pose datasets exist.
\end{itemize}
These factors make the problem significantly more challenging than standard 3D HPE from fixed ground cameras.

